% Circuit: GQSP layer matching qsp-subsection-v7.md (generic version).
% Replaces drafts/circuits/gqsp-thm6-laurent.tex and the current figure at
% 1_preliminaries.tex:1071--1113. Notation is aligned with v7 body:
%
%   * Generic Hermitian operator A with block encoding U_A and
%     subnormalisation alpha (no B, no b_pm, no alg:coh-specific objects).
%   * Walk operator W as a single box spanning the PREP ancilla and S
%     (the internal PREP-sandwich + SELECT structure is explained in v7's
%     walk-operator paragraph, Eq. eq:qsp-walk, and is NOT unrolled here).
%   * d (not d_JA) is the Jacobi-Anger truncation index. The circuit has
%     2d controlled signal slots: d controlled-W calls followed by d
%     controlled-W^dagger calls per MW Eq. 46 canonical "last k" pattern
%     (deg P = 2d, k = d).
%   * Output on S: L_d(W)|psi> + O(eps_QSP) = e^{-i delta A}|psi> + O(eps_QSP).
%
% Two dashed gategroups cover the d controlled-W slots and the d
% controlled-W^dagger slots. Rotation indices run j = 1, ..., d in the
% first group and j = d+1, ..., 2d in the second, consistent with the
% 2d+1 rotation triples produced by MW Alg. 1 on the underlying ordinary
% polynomial P(z) = z^d L_d(z) of degree 2d.
%
% Control convention: open (white) controls match the existing thesis
% figure at 1_preliminaries.tex:1071-1113 and the v7 parenthetical
% "(open-)controlled calls to W". Filled controls correspond to MW's
% canonical signal operator A = |0><0| x I + |1><1| x W and are a
% one-line edit if alignment with MW's definition is preferred.
%
\documentclass{article}
\usepackage[margin=1cm]{geometry}
\usepackage{tikz}
\usetikzlibrary{quantikz2}
\usepackage{amsmath,amssymb}
\usepackage{braket}

\newcommand{\ee}{\mathrm{e}}
\newcommand{\ii}{\mathrm{i}}

\tikzset{operator/.append style={rounded corners}}

\begin{document}
\pagestyle{empty}

\begin{figure}[ht]
  \centering
  \begin{quantikz}[row sep=0.7cm, column sep=0.55cm]
  %
  % Row 1: QSP ancilla. Opening R_0 (three parameters); within each
  % dashed group one (open-)control + one two-parameter rotation, repeated
  % d times. Rotation indices j = 1,...,d in the first group and
  % j = d+1,...,2d in the second.
  %
    \lstick{$\ket{0}_\mathrm{QSP}$}
      & \gate{R(\theta_0,\phi_0,\lambda_0)}
      & \octrl{1}
      \gategroup[3,steps=2,style={dashed,rounded corners,fill=black!6,inner sep=5pt},background,label style={label position=below,yshift=-0.45cm,align=center}]{\scriptsize Repeat $d$ times\\\scriptsize $j = 1,\dots,d$}
      & \gate{R(\theta_j,\phi_j,0)}
      & \octrl{1}
      \gategroup[3,steps=2,style={dashed,rounded corners,fill=black!6,inner sep=5pt},background,label style={label position=below,yshift=-0.45cm,align=center}]{\scriptsize Repeat $d$ times\\\scriptsize $j = d{+}1,\dots,2d$}
      & \gate{R(\theta_j,\phi_j,0)}
      & \rstick{$\bra{0}_\mathrm{QSP}$}\qw
      \\
  %
  % Row 2: PREP ancilla register b. Prepared in PREP|0>_b before the layer
  % and post-selected on PREP|0>_b afterwards. W acts on b+S jointly
  % (the PREP-sandwiched walk of MW Cor. 8, Eq. eq:qsp-walk).
  %
    \lstick{$\text{PREP}\,\ket{0}_b$}
      & \qwbundle{}
      & \gate[2]{W}
      & \qw
      & \gate[2]{W^\dagger}
      & \qw
      & \rstick{$\bra{0}\,\text{PREP}^\dagger$}\qw
      \\
  %
  % Row 3: System register S. Passes through the W / W^dagger slots; after
  % post-selection the top-left block on each qubitization subspace is
  % L_d(W) = e^{-i delta alpha cos theta} I + O(eps_QSP), i.e.
  % e^{-i delta A} + O(eps_QSP) on S (v7 body, target-polynomial paragraph).
  %
    \lstick{$\ket{\psi}$}
      & \qwbundle{}
      &
      & \qw
      &
      & \qw
      & \rstick{$L_d(W)\ket{\psi} + \mathcal{O}(\varepsilon_\mathrm{QSP})$}\qw
  \end{quantikz}
  \caption{GQSP layer of depth $\deg P = 2d$. With $k = d$ of the $2d$
  controlled signal slots deployed as controlled-$W^\dagger$
  (Thm.~6 of~\cite{motlagh2024generalized}, Eq.~46 canonical
  ``last $k$'' pattern) and the remaining $d$ as controlled-$W$, the
  walk $W$ is the PREP-sandwiched operator of (Eq.~\ref{eq:qsp-walk}).
  Post-selecting both ancillas on $\ket{0}_\mathrm{QSP}\otimes\text{PREP}\ket{0}_b$
  yields the Jacobi--Anger Laurent polynomial
  $L_d(W) = \ee^{-\ii\delta\alpha\cos\theta}\,I + \mathcal{O}(\varepsilon_\mathrm{QSP})$
  on each qubitization subspace, equivalently
  $\ee^{-\ii\delta A} + \mathcal{O}(\varepsilon_\mathrm{QSP})$ on the system
  register. The $2d+1$ rotation triples are extracted by MW Alg.~1 from
  the underlying ordinary polynomial $P(z) = z^{d} L_d(z)$ of degree $2d$.}
  \label{circ:gqsp}
\end{figure}

\end{document}
